% ************************** Thesis Abstract *****************************
% Use `abstract' as an option in the document class to print only the titlepage and the abstract.
\begin{abstract}

Docked bike-sharing systems (BSSs) are an accessible, low-carbon mobility option that can replace short urban car trips and improve connectivity to public transport. Over the past decade, BSSs have evolved into large-scale, data-rich services whose performance depends on decisions made across the full system life cycle, from station network design to day-to-day fleet operations. This thesis develops data-driven methods to support these planning and operational tasks, using Barcelona and its docked BSS, \textit{Bicing}, as real-world case studies.

To address the strategic design of bike-sharing networks, this thesis introduces a flexible framework for planning station locations. The framework combines modular multi-criteria scoring with network-based accessibility metrics that account for both the street network and local topography. By applying a genetic algorithm, it efficiently searches for optimal station layouts, enabling exploration of trade-offs between planning objectives and offering explicit control over features such as inter-station proximity and system accessibility. Applied to Barcelona, the framework identified near-optimal station configurations for different goals, including maximizing demand coverage, improving first- and last-mile connectivity, and promoting social equity. Importantly, the inclusion of altitude-adjusted metrics proves critical for equitable coverage, correcting for service gaps in hilly districts. Finally, an expansion for the l'Eixample district is presented, demonstrating how the framework can be applied to optimize new station locations to complement an existing network.

Turning to day-to-day operations, the thesis begins by characterizing the system's mobility patterns, demonstrating that electric bikes (e-bikes) are preferentially used for longer and uphill journeys compared to mechanical bikes (m-bikes). This insight underpins two predictive use cases. The first introduces a joint forecasting framework for e-bike fleet mobility and battery state by coupling a Markov-chain mobility model with a physics-based battery dynamics model. By simulating individual bike trajectories, this method generates system-wide forecasts. Validation results confirm the accuracy of the battery component, which estimates energy states with an error of less than 2\% ($RMSE=0.018$). Furthermore, the integrated framework achieves a notable correlation with observed station-level bike availability and average battery levels ($r \approx 0.70$ for both at a 30-minute horizon), providing reliable prediction windows that enable operators to shift from reactive retrieval to proactive charging strategies. 
    
The second use case addresses bike component reliability in both m-bikes and e-bikes through predictive maintenance using survival analysis and machine learning techniques. Combining historical mobility indicators, weather variables, and detailed maintenance orders, the models estimate failure risks for key components (brake pads, wheel spokes, and chains), achieving high predictive accuracy in all cases ($R^2 > 0.92$). Moreover, by leveraging interpretability techniques, the framework quantifies feature impacts, revealing that higher average speeds and electric propulsion significantly shorten the lifespan of these components, with usage intensity being the primary driver of degradation.

Overall, the thesis demonstrates that integrating computational methods into docked BSS planning and operations is feasible despite data limitations such as the absence of exact route information or high-resolution battery data. By enabling a transition from heuristic network design and reactive operations to data-driven planning and anticipatory strategies, the proposed frameworks provide scalable tools that can significantly improve system efficiency and sustainability. Furthermore, these contributions are transferable to other cities where comparable spatial and operational data are available, facilitating a shift towards rigorous, evidence-based urban mobility management.

\textbf{Keywords:} bike-sharing, electric mobility, spatial optimization, data-driven methods, station placement, battery prediction, predictive maintenance, Barcelona, Bicing.

\end{abstract}
